\subsection{Actual polynomial projective base change}
\label{subsec:native-projective-base-change}

The coefficient homomorphism is arbitrary. Every scheme and chart below is
constructed from the actual total-degree polynomial Proj and its actual
structure sheaf. The proofs retain all polynomial and section elements,
including nilpotents, without an independently supplied cartesian-square
or chart certificate.

\begin{theorem}[Whole coefficient and coordinate section maps]
\label{projective:actual-coefficient-map}
\lean{EnvelopeProperty.ProjBerk.coefficientGradedMap,
  EnvelopeProperty.ProjBerk.coefficientGradedMap_apply,
  EnvelopeProperty.ProjBerk.coefficientGradedMap_irrelevant_le,
  EnvelopeProperty.ProjBerk.nativeCoefficientMap,
  EnvelopeProperty.ProjBerk.nativeCoefficientMap_preimage_coordinateOpen,
  EnvelopeProperty.ProjBerk.homogeneousCoordinateMap,
  EnvelopeProperty.ProjBerk.nativeCoefficientMap_homogeneousChart,
  EnvelopeProperty.ProjBerk.coordinateDehom_map,
  EnvelopeProperty.ProjBerk.homogeneousCoordinateMap_dehom,
  EnvelopeProperty.ProjBerk.homogeneousCoordinateMap_equiv,
  EnvelopeProperty.ProjBerk.nativeCoefficientMap_coordinateChart,
  EnvelopeProperty.ProjBerk.nativeCoefficientMap_coordinateRestriction,
  EnvelopeProperty.ProjBerk.nativeCoefficientMap_coordinateSections,
  EnvelopeProperty.ProjBerk.isPullback_nativeCoefficientMap_coordinateChart}
\leanok
\uses{projective:native-structure-morphism,projective:native-relative-coordinate-chart,projective:native-scheme-coordinate-sections}
For any ring homomorphism $\varphi:R\to S$, the genuine graded coefficient
map induces an actual scheme morphism $g:\mathbb P^n_S\to\mathbb P^n_R$.
Its preimage of each coordinate open is the same coordinate open over $S$.
On that open its entire section map is coefficientwise polynomial mapping
$R[U_i:i\ne j]\to S[U_i:i\ne j]$ under the genuine section-ring
isomorphisms. Its square with the two actual coordinate inclusions is a
pullback.
\end{theorem}
\begin{proof}
Mapping coefficients preserves every homogeneous degree, even when some
coefficients become zero. Every target variable is the image of its original
variable, so the target irrelevant ideal is contained in the image of the
original irrelevant ideal. Apply Mathlib's actual graded Proj functor.
Its basic-open preimage formula gives the stated coordinate opens.
The actual homogeneous-localization chart square maps the whole numerator
and denominator of every fraction. Dehomogenizing at $X_j=1$ commutes
with coefficient mapping on every polynomial, by polynomial induction on
constants, addition and multiplication by a variable. Every actual
homogeneous fraction has such a numerator, so the genuine homogeneous
chart map is conjugate to the entire affine coefficient polynomial map.
Apply Spec and the actual chart isomorphisms to obtain the whole scheme
coordinate square. Restricting through the actual coordinate open immersions
gives the corresponding chart restriction. Its global-section functor map,
the real restriction top isomorphisms and the naturality of the Spec/global
section isomorphism give the stated equality on every section. Finally the
proved whole open-image preimage equality and the actual open-immersion
pullback criterion prove that the chart inclusion square is cartesian.
\end{proof}

\begin{theorem}[Actual polynomial Proj scalar extension]
\label{projective:actual-polynomial-proj-base-change}
\lean{EnvelopeProperty.ProjBerk.nativeCoordinateCover,
  EnvelopeProperty.ProjBerk.isPullback_coordinatePolynomialBaseChange,
  EnvelopeProperty.ProjBerk.isPullback_nativeCoefficientBaseChange,
  EnvelopeProperty.ProjBerk.nativeCoefficientMap_over_base,
  EnvelopeProperty.ProjBerk.nativeCoefficientBaseChangeIso,
  EnvelopeProperty.ProjBerk.nativeCoefficientBaseChangeIso_hom_fst,
  EnvelopeProperty.ProjBerk.nativeCoefficientBaseChangeIso_hom_snd}
\leanok
\uses{projective:actual-coefficient-map,projective:native-structure-morphism}
The whole coefficient square formed by $g$ and the two original coefficient
structure morphisms is cartesian. Thus there is a genuine scheme
isomorphism
\[
 \mathbb P^n_S\simeq
 \mathbb P^n_R\times_{\operatorname{Spec}R}\operatorname{Spec}S,
\]
whose first projection is $g$ and whose second projection is the actual
original $S$-structure morphism.
\end{theorem}
\begin{proof}
The coordinate open scheme inclusions form a real open cover: a relevant
homogeneous prime cannot contain every variable, since the genuine
irrelevant ideal is the ideal generated by these variables. On every
coordinate chart the whole ring square is the actual pushout
\[
 S\otimes_R R[U_i:i\ne j]\simeq S[U_i:i\ne j],
\]
with all original polynomials and both coefficient maps preserved. The
actual polynomial tensor algebra pushout, followed by contravariant Spec,
therefore gives a genuine affine pullback square. Conjugate by the true
pullback isomorphism of the preceding coordinate-inclusion square. The
actual coordinate structure identities then identify it with the restriction
of the proposed global coefficient square. Scheme pullback locality for this
actual open cover proves that the whole square is cartesian. Its whole
commuting coefficient diagram and the two projection identities follow from
the actual cartesian square and its pullback isomorphism. This construction
allows noninjective coefficient maps and nonreduced or zero rings; no field,
Noetherianity, valuation or finite-generation assumption on the base rings is
used.
\end{proof}

\begin{theorem}[Actual polynomial-coefficient model product]
\label{projective:actual-polynomial-line-product}
\lean{EnvelopeProperty.ProjBerk.nativePolynomialLineProductIso}
\leanok
\uses{projective:actual-polynomial-proj-base-change}
Specializing the genuine coefficient map to $R\to R[t]$ gives a scheme
isomorphism
\[
 \mathbb P^n_{R[t]}\simeq
 \mathbb P^n_R\times_{\operatorname{Spec}R}\operatorname{Spec}R[t].
\]
In particular the actual $k[t]$ projective model used by the checked
polynomial-Gauss ideal-log representation is the genuine algebraic product
presentation $\mathbb P^n_k\times\mathbb A^1_k$ used by the paper.
\end{theorem}
\begin{proof}
Take $S=R[t]$ and $\varphi=C$ in the actual scalar-extension isomorphism.
Both projections and the whole affine section maps are the same preserved
maps from the preceding producers. The existing Gauss bridge is therefore
on an actual scheme having this product presentation. This statement is an
algebraic scheme comparison and does not identify an analytic coherent
structure sheaf or claim a general scheme-product formula for arbitrary
analytic spaces.
\end{proof}
